\section{結び}
本論文において我々は, ストーリークリエーターの創作を支援する web アプリケーションを提案した. 本アプリケーションは, 物語構造分析に基づく LLM の活用を, 直感的な操作により実現するとともに, LLM とのインタラクティブなプロット創作手法を提供する. 商業誌に掲載される作品のプロット制作に利用してもらい, 本アプリケーションが高い操作性と満足のいく指示の反映をもたらし, プロットの質向上を目的とする LLM との共創手法に有用である可能性が示唆された. 今後は使用者の範囲を拡大した実証実験を実施し, より信頼性のある有用性の確認を行う. 

本研究がストーリー創作支援におけるクリエーターと LLM との共創のあり方に貢献することを期待する.
%私たちは、BunChoという日本語版GPT-2を使い、キーワードをタイピングすることで小説のタイトルやあらすじを生成できる日本の小説共創AIインターフェースを提案します。BunChoは利用者がTRPGに似た対話形式で創造的な小説を共創するのを支援します。私たちのユーザースタディでは、作家自身による書かれたあらすじとBunChoにより共創されたあらすじを比較しました。その結果、利用者は自分自身だけで作成するよりもBunChoと共創することをより楽しんでいることが示されました。私たちは、AIとの共創による日本の物語作りの未来の一つの方向性を示したと信じています。さらに、私たちはBunChoの共創におけるAI TRPGの効果とユーザーエクスペリエンスを評価する予定です。